\begin{abstract}
Test-time compute scaling---repeated sampling with self-consistency or best-of-$N$---yields large gains on discrete-answer reasoning, but whether these gains transfer to NER remains open.
We characterize when scaling succeeds and fails, using NER (four datasets, primarily 7--14B fine-tuned) as the primary analysis, with preliminary cross-task pilots (RE, SQL).
Quality estimation signals predict output quality (AUROC up to $0.82$), yet instance-level selection yields no detectable gains (${<}1$\,pp; CI crosses zero), revealing a \textbf{correlation-selection gap}.
A diagnostic framework identifies two co-occurring conditions---\textit{degeneracy} (i.e., multiple outputs collapsing to identical predictions) and signal alignment collapse---that are jointly associated with scaling returns; cross-task exploration confirms task-specific failure modes.
Entity-level construction yields conditional gains ($+1.40$\,pp on Few-NERD under 3-epoch training; collapses at convergence; global Holm $p_{\text{adj}}{=}0.055$, suggestive; per-family $p_{\text{adj}}{=}0.006$, post-hoc family partition) but exhibits inverse scaling in most configurations.
The framework is primarily diagnostic, with limited predictive power for identifying favorable conditions, informing resource allocation decisions before investing in repeated inference.
\end{abstract}
